% Reviewer 1
\reviewer
\medskip
\textcolor{blue}{We sincerely thank Reviewer 1 for the confirmation that several comments from the previous round have been successfully addressed. We appreciate the opportunity to further refine our work and enhance its quality. We have substantially revised the manuscript to address the remaining comments, with detailed point-by-point responses provided below.} 

\noindent \textit{Note: Reviewer 1 in this round was Reviewer 2 in the previous round.}
\medskip

\begin{revcomment}[1.0]
	Thanks for the effort to revise the paper. After the revision, the paper has improved significantly. However, there are still some problems raised after the revision that I hope the authors can resolve. 
\end{revcomment}

\hypertarget{Response to R1}{}
\begin{revresponse}[1.0]
    Thank you for the time and effort invested in this review. We appreciate the positive feedback regarding the improvements in our manuscript. Below, we address your remaining concerns point by point. 
\end{revresponse}

\begin{revcomment}[1.1]
	My first question is about the domain knowledge in transportation. I think the developed theory might be too generalized. I mean if this paper is submitted to an OR or management science journal, there will be no problem at all. But it is for Transportation Science. In order to fall in the scope of TS, it should have transportation specific domain application. 
    
    \textbf{(1.1.1)} The author replied that ``crowd carriers' route choice, platform's matching/pricing strategy, and the interplay between crowdshipping and dedicated delivery'' are fundamental characteristics in crowdshipping. However, almost all mechanism design problems have matching/pricing strategy and the interplay between agents, which do not make the problem specific to crowdshipping. The only element specific to crowdshipping is the route-choice-based bidding. However, the model (The Winner Determination Problem) does not model the route-choice-based bidding. What all matters is the bundle $s$ in the model. The relationship between bundle and route choice is not clarified. 
    
    \quad Is it possible to model the bidding for bundle based on the route choice? The current model looks like to enumerate all possible bundles, each of which corresponds to a route. In another words, all routes are predetermined, so that the platform simply recommends routes to agents for bidding. Is it possible to bid like this way: let the service provider bid for a cost rate per mile for delivery? And model it into the formulas?
    
    \textbf{(1.1.2)} Then, another question. The developed model might also be applied in shared mobility, such as carpooling, right? Any difference for the modeling? 
\end{revcomment}

\begin{revresponse}[1.1.1]
    Thanks for this thoughtful comment. Our fundamental motivation stems directly from the transportation domain, specifically the emerging paradigm of crowdshipping. We observed that crowdshipping systems exhibit unique, previously uncharacterized features, that is, (1) the original trip purpose and route choice of crowd carriers and (2) the interplay between crowdshipping and dedicated delivery services. We identified combinatorial auctions utilizing XOR bids as the ideal theoretical tool to model and analyze the distinct features of crowdshipping. The relationship between route choice and bids is completely captured before the WDP is formulated:
    
    \begin{itemize}
        \item Phase 1: Route-Driven Bundle Generation (The Transportation Foundation). A bundle is not an abstract combination of orders; it is strictly derived from a carrier’s physical routing reality and human constraints (e.g., origin, destination, detour tolerance). The system calculates the optimal detour route for any feasible bundle and translates this travel cost directly into the bid. Thus, route choices are completely embedded into the bids prior to the auction.
        \item Phase 2: The WDP (The Allocation Mechanism). The WDP takes these route-encoded bids from Phase 1 as fixed inputs. It determines the winner and finds the system-optimal assignment by evaluating these crowdsourced bids against dedicated backup services.
    \end{itemize}

    Because the bids already encapsulate the carriers' routing decisions, the WDP does not need to model route choices endogenously. To make this architecture transparent, we have substantially revised \added{Section 3.2} to explicitly detail how physical routes are generated and translated into bids. 
    
    One way to structure bids is to let carriers submit a cost rate per mile, which can then be incorporated into the model. However, this approach assumes that a carrier's preferences can be captured by a small set of parameters, which may not hold given the heterogeneity in how detour costs are formed across individuals.
    
    Parameterized bidding is often used when the bid space becomes complex and the number of feasible bundles is large. In our crowdshipping setting, however, carriers are occasional participants with tight constraints: small capacity and limited detour tolerance keep the number of feasible bundles small and manageable, making such parameterization unnecessary. 
    
    \begin{changes}
        (Section 3.2) xxx
    \end{changes}
\end{revresponse}


\begin{revresponse}[1.1.2]
    We greatly appreciate this insightful and forward-looking question. To answer directly: Yes, our model can indeed be naturally extended and applied to shared mobility systems such as carpooling and ride-pooling. 
    
    The core structure of our framework — including occasional carriers with pre-existing trips, bundle feasibility derived from route optimization, detour cost-based XOR bids, and a backup service provider as the cost benchmark — is generic to a broad family of pickup-and-delivery shared-use systems, all of which fall into the General Pickup and Delivery Problem (GPDP) framework (Savelsbergh and Sol, 1995). 

    At the same time, we fully agree that certain modeling differences must be carefully addressed when adapting our model to carpooling:
    \begin{itemize}
        \item Market structure: Our current model is designed for a one-sided reverse combinatorial auction, where the platform acts as the sole buyer of delivery services, and we only need to encourage incentive compatibility on the supply side (crowd carriers). By contrast, carpooling is typically a two-sided matching market involving both drivers and riders, which requires additional consideration of riders' private information, participation constraints, and payment preferences.
        \item Routing problem type: Our delivery model corresponds to a Pickup and Delivery Problem (PDP), where constraints focus mainly on the courier’s global detour limit and task time windows. Carpooling belongs to the Dial-a-Ride Problem (DARP), a stricter subclass of GPDP, which imposes additional individual-level constraints such as each rider’s maximum allowed travel time, in-vehicle waiting time, and dynamic passenger capacity limits along the route.
        \item Capacity constraint: In our delivery setting, capacity is a static limit on the number of orders. In carpooling, capacity is a dynamic constraint on the number of passengers in the vehicle at any time, which changes en route and further restricts feasible
        bundles.
        \item Backup option interpretation: In our model, the backup (third-party logistics) serves as a platform-level fulfillment guarantee to ensure all orders are delivered. In carpooling, the backup (e.g., taxi or ride-hailing service) acts as an individual outside option for each rider, rather than a centralized sourcing contract. 
    \end{itemize}
    We have emphasized this generalizability and its potential for shared mobility in the revised manuscript (\added{Section 1.1}). Following your suggestion, we will continue to explore this promising direction in our future work. 
    \begin{changes}
        (Section 1.1) To our knowledge, this framework is among the first to propose analytically tractable and scalable auction mechanisms for hybrid crowdshipping systems that accurately capture crowd carriers' preferences. It offers a practical approach to order allocation and payment design, effectively incentivizing heterogeneous travelers to serve as crowd carriers, and can be readily adapted to other shared mobility systems such as carpooling. More broadly, this work advances combinatorial auction theory with applicability to a wide range of structurally analogous problems beyond transportation, that is, settings where service providers bid on mutually exclusive bundles of items and each item has a fixed-price backup option. 
    \end{changes}  
\end{revresponse}


\begin{revcomment}[1.2]
	My second question is about one-to-one vs many-to-many. Based on the authors' response, I guess that one-to-one means the following: If an agent wins A, his value is 10 dollars. If the agent wins B, his value is 15 dollars. Then if he wins both, then his value is 10+15 = 25 dollars. Many-to-many means if agent wins both, the value is not necessarily 25 dollars. Is my understanding right? 
    
    \quad If yes, it is incorrect that the author claims ``The auctions may not fully capture the non-additive valuation of order bundling by crowd carriers and often assign orders one by one to crowd carriers.'' Almost all combinatorial auction models do not rely on the assumption of ``additive valuation''. This is not a novel contribution of this paper. 
\end{revcomment}

\begin{revresponse}[1.2]
	Thank you very much for this thoughtful comment. Your understanding of one-to-one vs. many-to-many and additive/non-additive valuations is right. One-to-one allocation implicitly assumes that the value (cost) of a bundle is the sum of values (costs) of individual orders, i.e., additive valuation. Many-to-many (bundle) allocation allows the value of a bundle to differ from the sum, i.e., non-additive (subadditive/superadditive) valuation. 

    In the literature, some studies explored auction mechanisms for crowdshipping that incorporate both crowd carriers and dedicated delivery services. They do not consider crowd carriers' original trip purposes and rely on methods with limited analytical tractability and scalability. Other studies achieved scalability but at the cost of imposing strong assumptions on crowd carriers' preference structures. For example, they assume that the cost of serving orders in a bundle can be decomposed as the sum of cost serving each order. 
    By contrast, our framework strikes a balance between both theoretical economic properties and computational efficiency for real-world large-scale applications, while retaining the full non-additive (specifically, subadditive) detour cost structure in a bundle.

    The quote you mentioned appeared in \added{Section 2.3} where we summarize the research gaps from the existing literature. We did not mention this as a novel contribution, it was merely a description of the issue. 
    After reviewing your comment, we realized it might be inaccurate, we have revised the ambiguous statement in the manuscript as follows:
    \begin{changes}
        (Section 2.3) To summarize, few existing auction mechanisms for crowdshipping have studied systems that accommodate both dedicated delivery services and crowd carriers, while balancing theoretical economic properties and computational scalability. Among the few that do, they impose strong assumptions about crowd carriers' preference structures. We set out to fill these gaps.
    \end{changes}
\end{revresponse}

\begin{revcomment}[1.3]
	I raised this question about fair comparison of the proposed mechanism by greedy approach with VCG mechanism. The author claimed that both the proposed mechanism and VCG mechanism targets at minimizing the total cost, which is the negative social welfare. This is understandable. 
    
    \textbf{(1.3.1)} However, when we look at the comparison in Figure 3, we found that the cost of VCG outperforms the greedy mechanism significantly in some cases, such as M =3, N =5: VCG cost is 735 while Greedy cost is 1250. This result is not satisfactory. Although, the author also compares the two mechanisms in terms of profit, fulfillment, utilization rate, etc., none of these metrics is the objective of the designed mechanism. Maybe the author should clarify the objective of the designed mechanism and select a good benchmark for fair comparison. 

    \textbf{(1.3.2)} Also, I realized that the authors proved a bound for social welfare, which is very impressive. But can you show the detailed values of these bounds for specific examples, say M =3, N =5: gap between VCG and Greedy is 735 -1250, then what is the upper bound of this gap? 
\end{revcomment}

\begin{revresponse}[1.3.1]
	Thank you very much for this thoughtful and rigorous comment. 
    
    We agree with your observation and concern: As the theoretically optimal mechanism that exactly solves the winner determination problem (WDP), VCG by definition achieves the minimal possible total cost, so it is natural and expected that VCG significantly outperforms the greedy mechanism in total cost in small-scale scenarios (such as M=3, N=5). The cost difference you pointed out reflects the efficiency loss from approximation, which is exactly the trade‑off we want to highlight. We now clarify the following key points explicitly and will revise the manuscript accordingly: 
    \begin{itemize}
        \item The unique, unified objective of both mechanisms. Both the VCG mechanism and the proposed greedy mechanism share the same core objective: to minimize the total fulfillment cost of the platform. This remains the only optimization objective of our mechanism design. All other metrics (carrier profit, fulfillment rate, utilization rate) are supplementary operational indicators, not optimization objectives, and we have clearly marked this distinction in the revised paper. 
        \item Role of VCG as the ideal theoretical benchmark. VCG is used as the optimal benchmark because it achieves the globally minimal total cost while satisfying incentive compatibility, individual rationality, and budget balance. It represents the performance upper bound that any feasible mechanism can possibly reach in this setting.
        \item Purpose of the greedy mechanism: computational tractability for large‑scale, real‑time systems. The greedy mechanism is proposed not to outperform VCG in total cost, but to resolve the critical limitation of VCG: it is computationally intractable for large‑scale real‑time crowdsourced logistics systems (where the number of orders and agents can be thousands). The greedy mechanism achieves exponentially faster computation with a controllable efficiency loss, making it practically implementable in real‑world applications.
        \item Fairness and meaning of the comparison. Our comparison is not intended to claim that the greedy mechanism is ``better'' than VCG. Instead, it illustrates the fundamental trade‑off in real‑time shared mobility systems: VCG achieves optimal cost, but computationally infeasible at scale; Greedy is near-optimal cost, while highly efficient, and implementable in practice. 
    \end{itemize}
    
    We have revised Figure 3 and its discussion to: (1) clearly label VCG as the theoretical optimal benchmark; (2) explicitly explain that the cost difference reflects the intentional trade-off between optimality and computational complexity; (3) emphasize that supplementary performance metrics are only for operational reference and not the design objectives of our mechanisms.   
    
    \begin{changes}
        xxx
    \end{changes} 
\end{revresponse}




\begin{revresponse}[1.3.2]
	Thank you very much for this insightful and rigorous comment, which helps us clearly explain the meaning of our theoretical worst-case upper bound, the gap between the bound and empirical performance, and the underlying driving factors in our experiments. We fully address your concerns as follows.
    
    First, the upper bound we derived in the theory is a universal worst-case guarantee that holds for all possible input instances, including extremely unfavorable scenarios where there is little synergy between orders, few feasible matches between crowd carriers and orders, or most tasks have to be fulfilled by the high-cost backup service. This bound provides a reliable performance limit that the greedy mechanism will never violate in practice. Since it is constructed for the most conservative and extreme cases, it is expected to appear loose relative to realistic, well-structured inputs — this is a standard and well-accepted property in approximation algorithm analysis.
    
    The theoretical worst-case upper bound in the small-scale case (M=3, N=5) is 4307, while the actual cost of the greedy mechanism is 1250. The large gap can be fully explained by the extremely high matching rate (around 0.9) of orders fulfilled by crowd carriers in this setting. Strong synergies exist between crowd carriers' direct routes and delivery orders, and the greedy mechanism effectively captures these bundle synergies. As a result, the actual efficiency loss between the greedy mechanism and the optimal VCG mechanism is naturally much smaller than the theoretical worst-case upper bound.
    
    We further validate this explanation using our experiments based on real-world mobility datasets from Atlanta and Fedex Sameday delivery rates. In these realistic instances, only about 30\% of orders can be served by crowd carriers, which is much closer to the conservative assumptions used in our worst-case bound derivation. In this case, the theoretical upper bound becomes considerably tighter and aligns much better with the optimal VCG performance. This contrast confirms that the tightness of our upper bound naturally depends on the structure of the input, especially the share of orders fulfilled by crowd carriers, while the theoretical guarantee itself remains universally valid.
    
    Notably, even when the worst-case bound appears loose, the actual total cost difference between the greedy mechanism and the optimal VCG mechanism remains small and well-controlled in all our experiments. The high theoretical upper bound does not imply poor practical performance. In addition, as the problem scale increases, although the theoretical worst-case gap may widen, the empirical performance of the greedy mechanism remains close to the optimal VCG outcome, because it can effectively capture the most valuable bundle synergies even in large-scale real-time systems. 
    
    Following your suggestion, we have revised the manuscript to: (1) clearly emphasize that our upper bound is a universal worst-case guarantee rather than a tight bound for realistic instances; (2) explicitly explain how the share of orders fulfilled by crowd carriers affects the tightness of the bound; (3 supplement the numerical comparison between the theoretical bound, the greedy cost, and the VCG cost for both synthetic and real-data instances; (4) clarify the efficiency performance of the greedy mechanism as the problem scale grows.

    We greatly appreciate your valuable feedback, which significantly improves the transparency and rigor of our theoretical and experimental analysis.
    \begin{changes}
        xxx
    \end{changes}
\end{revresponse}


\begin{revcomment}[1.4]
	This is minor issue. I suggest putting the numbers in tables so that the comparison is straightforward. 
\end{revcomment}

\begin{revresponse}[1.4]
	We have adopted your suggestion and formalized the comparative numerical parameters directly into concise tables, greatly improving spatial readability.
    \begin{changes}
        xxx
    \end{changes}
\end{revresponse}
